%! Unit 2: Functions, Transformations, and Graph Analysis — Summative Assessment — Practice Test --- Study Copy (blank)
%
% Blueprint (100 pts): Part A vocabulary 8x1, Part B multiple choice 6x2,
% Part C short answer 10x5 (ONE ITEM PER CONTENT LESSON, in lesson order, then
% three synthesis items), Part D extended response 2x15. The actual test is the
% parallel form: same blueprint, different numbers, reshuffled vocabulary and
% multiple-choice letters.
%
% UNIT 2 IS A GRAPH-READING UNIT, so unlike Unit 1's test this one carries five
% pre-drawn pgfplots figures (C4, C5, C6, C9, D1). That is the AFDA guidance's own
% scope — characteristics are investigated from a graph — and it is why this form
% runs longer than Unit 1's four pages. No item asks a student to sketch.
%
% WORK RULE: every computation is followed by a \begin{work} block authored
% byte-identically here and in ../../test_keys/practice_test_key/main.tex. Under
% -boxes the block ships as a \vphantom, so the student gets exactly the room the
% answer needs; under -key the same block prints in keyred. That is what keeps
% this test and its key the same number of pages — which matters because
% shared/unit.mk swaps sample_test_key in for sample_test at the tail of the key
% packet. \workrowsep is the handwriting room, and it moves both files together.
\documentclass[10pt]{article}
\usepackage{atda-article}
\usepackage{atda-boxes}
\setlength{\workrowsep}{6pt}

% Part-divider strip.
% BOXGUARD: the guard lives inside the macro, so a part header can never be
% stranded alone at the foot of a page, in the blank or the key — both files use
% this identical definition.
\newcommand{\parthead}[1]{%
  \par\boxguard[9]%
  \vspace{6pt}\noindent
  \begin{headlinebox}{royal}{\color{white}\bfseries #1}\end{headlinebox}%
  \vspace{4pt}}

% Shared axis styling for the five figures. Defined identically in all four files.
\pgfplotsset{testaxis/.style={
    grid=both, grid style={linegray!45}, minor grid style={linegray!30},
    axis lines=left, tick label style={font=\tiny}, label style={font=\scriptsize}}}

\begin{document}

\pageheader{Unit 2: Functions, Transformations, and Graph Analysis}{Practice Test --- Study Copy}
\namedateperiod

\vspace{0.10in}
\boxguard[8]
\begin{remindbox}
\small
\textbf{This is a practice test.} It mirrors the real Unit 2 test in length, format,
and the ideas it covers, but the numbers and contexts differ. Work every problem
with no notes, then check your answers against the key.
\end{remindbox}

% ======================================================================
\parthead{Part A --- Vocabulary \quad (8 items, 1 point each --- 8 points)}

\noindent\small
Write the letter of the best definition on the line beside each term. \textbf{Two definitions
will not be used.}

\vspace{5pt}
\noindent
\begin{tabularx}{\linewidth}{@{} p{0.95cm} X p{0.95cm} X @{}}
\blank{0.7cm} & \textbf{1.} Domain
  & \blank{0.7cm} & \textbf{5.} Zero of a function \\[4pt]
\blank{0.7cm} & \textbf{2.} Range
  & \blank{0.7cm} & \textbf{6.} Absolute maximum \\[4pt]
\blank{0.7cm} & \textbf{3.} Parent function
  & \blank{0.7cm} & \textbf{7.} Horizontal asymptote \\[4pt]
\blank{0.7cm} & \textbf{4.} Translation
  & \blank{0.7cm} & \textbf{8.} Piecewise-defined function \\
\end{tabularx}

\vspace{6pt}
\begin{enumerate}[label=(\Alph*), leftmargin=2.4em, itemsep=1pt, topsep=2pt, font=\bfseries]
  \small
  \item The set of all output values a function produces.
  \item A horizontal line that a graph approaches at one or both ends without ever reaching it.
  \item The simplest function in a family, before any shift, flip, stretch, or compression.
  \item A stretch or a compression --- a change to the \emph{shape} of a graph.
  \item The set of all input values a function is allowed to take.
  \item A function defined by different rules on different stretches of its domain.
  \item The point where a graph crosses the vertical axis.
  \item A slide that moves a graph up, down, left, or right without changing its shape.
  \item An input value for which the output is $0$.
  \item The largest output value a function reaches over its \emph{entire} domain.
\end{enumerate}

% ======================================================================
\parthead{Part B --- Multiple Choice \quad (6 items, 2 points each --- 12 points)}

\noindent\small Circle the best answer.

\begin{enumerate}[leftmargin=1.8em, itemsep=6pt, topsep=5pt]
\small

\item The point $(3,10)$ is on the graph of $f$. Which statement says the same thing?
\begin{enumerate}[label=(\Alph*), leftmargin=2.2em, itemsep=1pt, topsep=2pt]
  \item $f(10)=3$
  \item $f(3)=10$
  \item $f(3)=3$
  \item $f(10)=10$
\end{enumerate}

\item The graph of $y=(x+5)^2$ is the graph of $y=x^2$ moved
\begin{enumerate}[label=(\Alph*), leftmargin=2.2em, itemsep=1pt, topsep=2pt]
  \item right $5$
  \item up $5$
  \item left $5$
  \item down $5$
\end{enumerate}

\item A yearbook's profit is \$0 at prices of \$0 and \$60, and rises to a high of \$1800 at a
price of \$30. Over which \textbf{prices} is the profit increasing?
\begin{enumerate}[label=(\Alph*), leftmargin=2.2em, itemsep=1pt, topsep=2pt]
  \item from $0$ to $1800$
  \item from $0$ to $60$
  \item from $30$ to $60$
  \item from $0$ to $30$
\end{enumerate}

\item Which statement about the graph of $y=2^x$ is true?
\begin{enumerate}[label=(\Alph*), leftmargin=2.2em, itemsep=1pt, topsep=2pt]
  \item It crosses the $x$-axis at $x=0$.
  \item It has no $x$-intercept: the curve approaches the line $y=0$ without reaching it.
  \item It has an $x$-intercept somewhere far to the left.
  \item It has neither an $x$-intercept nor a $y$-intercept.
\end{enumerate}

\item A garden center pays \$12 an hour for the first $40$ hours of a week and \$18 an hour after
that. The graph of a week's pay $P(h)$ has a \textbf{corner} at $h=40$, where the rule changes.
Which statement is true?
\begin{enumerate}[label=(\Alph*), leftmargin=2.2em, itemsep=1pt, topsep=2pt]
  \item The graph has a break at $h=40$, so $P$ is discontinuous there.
  \item $P(40)$ has two different values, one from each rule.
  \item $P$ is continuous at $h=40$: both rules give the same pay there, so the graph can be drawn
        without lifting your pencil.
  \item $h=40$ is not in the domain of $P$.
\end{enumerate}

\item A set of data curves \textbf{upward} as $x$ increases. What does that observation tell you
about which family fits?
\begin{enumerate}[label=(\Alph*), leftmargin=2.2em, itemsep=1pt, topsep=2pt]
  \item The data must be exponential.
  \item The data must be quadratic.
  \item Nothing at all.
  \item It rules out linear --- and nothing more. You still have to check the differences and the
        ratios.
\end{enumerate}

\end{enumerate}

% ======================================================================
\parthead{Part C --- Short Answer \& Computation \quad (10 items, 5 points each --- 50 points)}

\noindent\small Show all work. An answer with no supporting work earns no credit.

\begin{enumerate}[leftmargin=1.8em, itemsep=4pt, topsep=5pt]
\small

\item A pool is being drained. After $t$ hours it holds $W(t) = 3000-250t$ gallons, for
$0 \le t \le 12$. \ (a) Find $W(4)$ and say what it means. \ (b) Solve $W(t)=750$ and say what it
means. \ (c) Write the domain and the range in interval notation.
\begin{work}
  \text{(a)} \quad W(4) &= 3000-250(4) = 2000 \text{ gallons after 4 hours} \\
  \text{(b)} \quad 3000-250t = 750 \;\Rightarrow\; 250t &= 2250 \;\Rightarrow\; t = 9 \text{ hours} \\
  \text{(c)} \quad \text{domain } [0,12], \quad \text{range } &[0,3000] \quad (W(12)=0,\; W(0)=3000)
\end{work}

\item For each equation, describe the transformation from its parent function and give the range.
\quad (a) $y=(x-4)^2+1$ \quad (b) $y=2^x-5$ \quad (c) $y=-x^2$
\begin{work}
  \text{(a)} \quad &x^2 \text{ shifted right } 4 \text{ and up } 1; \quad \text{range } [1,\infty) \\
  \text{(b)} \quad &2^x \text{ shifted down } 5; \quad \text{range } (-5,\infty) \\
  \text{(c)} \quad &x^2 \text{ reflected over the } x\text{-axis}; \quad \text{range } (-\infty,0]
\end{work}

\item A parabola has vertex $(2,-3)$ and passes through the point $(4,5)$. \ (a) Write its
equation in the form $y=a(x-h)^2+k$. \ (b) Verify your equation by substituting the point
$(1,-1)$.
\begin{work}
  \text{(a)} \quad 5 &= a(4-2)^2 - 3 \\
  8 &= 4a \;\Rightarrow\; a = 2, \qquad y = 2(x-2)^2-3 \\
  \text{(b)} \quad 2(1-2)^2-3 &= 2(1)-3 = -1 \quad \checkmark
\end{work}

% BOXGUARD: without this, item 4's stem — the two lines that define what A(t) MEANS
% and that A<0 is below the platform — sat alone at the foot of p2, with the figure
% and all four questions opening p3. That is a student reading the context on one
% page and answering on the next, and it is the same defect class the lessons hit at
% 2.4 item 4, 2.5 item 8 and 2.6 item 4. Invisible to page counts: the key stranded
% the identical two lines, so parity was perfect at 6/6. [16] is sized to the stem
% plus the 5.2cm axis, so the two can never separate. Mirrored in the key.
\boxguard[16]
\item A drone flies out over a canyon. $A(t)$ is its altitude in meters \textbf{relative to the
launch platform} ($A<0$ means below the platform), $t$ minutes into the flight.

\begin{center}
\begin{tikzpicture}
\begin{axis}[testaxis, width=11.0cm, height=5.2cm,
    xlabel={$t$ (minutes)}, ylabel={$A$ (meters)},
    xmin=0, xmax=12, ymin=-7.4, ymax=7.4,
    xtick={0,2,4,6,8,10,12}, ytick={-6,-4,-2,0,2,4,6},
    minor x tick num=1, minor y tick num=1]
  \addplot[royal, very thick] coordinates {(0,6) (2,0) (5,-6) (8,0) (10,4) (12,4)};
\end{axis}
\end{tikzpicture}
\end{center}

\vspace{-2pt}
(a) Give the $A$-intercept and say what it means. \ (b) Give both zeros and say what a zero means
here. \ (c) Give the intervals on which $A$ is increasing, decreasing, and constant. \ (d) Give
the absolute maximum and the absolute minimum --- each as a \textbf{value and a location}.
\begin{work}
  \text{(a)} \quad &(0,6): \text{ the drone starts } 6 \text{ m above the platform} \\
  \text{(b)} \quad &t=2 \text{ and } t=8: \text{ the drone is level with the platform} \\
  \text{(c)} \quad &\text{decreasing } (0,5), \quad \text{increasing } (5,10), \quad
      \text{constant } (10,12) \\
  \text{(d)} \quad &\text{maximum } 6 \text{ m at } t=0; \qquad
      \text{minimum } -6 \text{ m at } t=5
\end{work}

\item After $w$ weeks of practice a typist's speed is $S(w) = 90-80(0.5)^w$ words per minute.

\begin{center}
\begin{tikzpicture}
\begin{axis}[testaxis, width=9.6cm, height=4.6cm,
    xlabel={$w$ (weeks)}, ylabel={$S$ (wpm)},
    xmin=0, xmax=6, ymin=0, ymax=100,
    xtick={0,1,2,3,4,5,6}, ytick={0,20,40,60,80,100}, minor y tick num=1]
  \addplot[slate, dashed, thick, domain=0:6, samples=2]{90};
  \addplot[royal, very thick, domain=0:6, samples=120]{90-80*0.5^x};
\end{axis}
\end{tikzpicture}
\end{center}

\vspace{-2pt}
(a) Write the equation of the horizontal asymptote. \ (b) Describe the end behavior as
$w \to \infty$. \ (c) Will this typist ever reach exactly $90$ wpm? \ (d) Could this typist ever
reach $95$ wpm? \textbf{Give a different reason for (c) and for (d).}
\begin{work}
  \text{(a)} \quad &y = 90 \\
  \text{(b)} \quad &\text{as } w \to \infty, \; S \text{ levels off toward } 90 \\
  \text{(c)} \quad &\text{No --- the speed gets closer and closer to } 90
      \text{ but never arrives.} \\
  \text{(d)} \quad &\text{No --- } 95 \text{ is on the far side of the asymptote entirely.}
\end{work}

\item A rideshare charges a flat \$6 for the first $3$ miles, then \$2 for each mile after that.
The fare for a ride of $m$ miles is $F(m)$.

\begin{center}
\begin{tikzpicture}
\begin{axis}[testaxis, width=9.6cm, height=4.6cm,
    xlabel={$m$ (miles)}, ylabel={$F$ (dollars)},
    xmin=0, xmax=10, ymin=0, ymax=22,
    xtick={0,1,2,3,4,5,6,7,8,9,10}, ytick={0,4,8,12,16,20}, minor y tick num=1]
  \addplot[royal, very thick] coordinates {(0,6) (3,6) (10,20)};
\end{axis}
\end{tikzpicture}
\end{center}

\vspace{-2pt}
(a) Find $F(2)$ and $F(7)$. \ (b) Name the boundary point. \ (c) Is $F$ continuous or
discontinuous at the boundary? \textbf{Justify from the graph.} \ (d) In one sentence, say what
the two rules mean to a rider.
\begin{work}
  \text{(a)} \quad F(2) &= 6 \text{ dollars}; \qquad F(7) = 6 + 2(7-3) = 14 \text{ dollars} \\
  \text{(b)} \quad &m = 3 \\
  \text{(c)} \quad &\text{Continuous: both rules give } \$6 \text{ at } m=3,
      \text{ so the pieces meet.} \\
  &\text{It is a corner, not a break --- no pencil comes off the paper.} \\
  \text{(d)} \quad &\text{The first 3 miles cost a flat } \$6;
      \text{ every mile after that adds } \$2.
\end{work}

\item For each table the $x$-values step by the same amount. Name the family --- \textbf{linear,
quadratic, or exponential} --- and justify with the differences or the ratio.

\vspace{3pt}
\noindent
\begin{tabularx}{\linewidth}{@{} X X X @{}}
\begin{tabular}{@{}c|cccc@{}}
\multicolumn{5}{@{}l}{\textbf{(a)}}\\
$x$ & 0 & 1 & 2 & 3 \\ \hline
$y$ & 4 & 7 & 10 & 13 \\
\end{tabular}
&
\begin{tabular}{@{}c|cccc@{}}
\multicolumn{5}{@{}l}{\textbf{(b)}}\\
$x$ & 0 & 1 & 2 & 3 \\ \hline
$y$ & 5 & 10 & 20 & 40 \\
\end{tabular}
&
\begin{tabular}{@{}c|cccc@{}}
\multicolumn{5}{@{}l}{\textbf{(c)}}\\
$x$ & 0 & 1 & 2 & 3 \\ \hline
$y$ & 2 & 5 & 14 & 29 \\
\end{tabular}
\end{tabularx}
\begin{work}
  \text{(a)} \quad &\text{first differences } 3,3,3 \text{ --- constant} \;\Rightarrow\;
      \textbf{linear} \\
  \text{(b)} \quad &\text{ratios } 2,2,2 \text{ --- constant} \;\Rightarrow\;
      \textbf{exponential} \\
  \text{(c)} \quad &\text{first differences } 3,9,15; \;
      \text{second differences } 6,6 \;\Rightarrow\; \textbf{quadratic}
\end{work}

\item (a) Write $-2 \le x < 7$ in interval notation, and say what each bracket or parenthesis
means. \ (b) Write ``all outputs greater than $3$'' in interval notation. \ (c) A field trip costs
\$8 per student and the bus holds at most $30$ students. Give the domain of $C(n)=8n$ in this
context, and explain why interval notation cannot express it.
\begin{work}
  \text{(a)} \quad &[-2,7): \text{ the bracket includes } -2, \text{ the parenthesis excludes } 7 \\
  \text{(b)} \quad &(3,\infty) \quad
      (\infty \text{ never takes a bracket}) \\
  \text{(c)} \quad &\{0,1,2,\dots,30\} \text{ --- only \emph{whole} students ride the bus, and an} \\
  &\text{interval would include values like } 17.4, \text{ which is not a domain value.}
\end{work}

\item A school video has $V(t) = 3 \cdot 2^{\,t}$ thousand views $t$ days after it is posted.

\begin{center}
\begin{tikzpicture}
\begin{axis}[testaxis, width=9.2cm, height=4.6cm,
    xlabel={$t$ (days)}, ylabel={$V$ (thousand views)},
    xmin=0, xmax=5, ymin=0, ymax=100,
    xtick={0,1,2,3,4,5}, ytick={0,12,24,36,48,60,72,84,96}, minor x tick num=1]
  \addplot[royal, very thick, domain=0:5, samples=120]{3*2^x};
\end{axis}
\end{tikzpicture}
\end{center}

\vspace{-2pt}
(a) Find $V(5)$ \textbf{algebraically}. \ (b) Use the \textbf{graph} to find the day the video
reaches $48$ thousand views. \ (c) Name which method you used for each, and say why (a) is the
easier one to trust here.
\begin{work}
  \text{(a)} \quad V(5) &= 3 \cdot 2^5 = 3 \cdot 32 = 96 \text{ thousand views} \\
  \text{(b)} \quad &\text{read across from } 48 \text{ to the curve, then down: } t = 4 \text{ days} \\
  \text{(c)} \quad &\text{(a) is algebraic --- substitution; (b) is graphical --- a read.} \\
  &\text{Substitution gives an exact value; a graph read is only as good as the grid.}
\end{work}

% BOXGUARD: without this, item 10's stem sat at the foot of p4 and its whole work
% block — five rows of answer room — opened p5 as an unlabelled blank strip above
% the Part D banner. A student would be writing five answers on a page that shows
% no question. Same defect class as item 4 above, and equally invisible to the page
% count (the key printed the same five rows in the same place, so parity stayed
% 6/6). [14] is the stem plus the five work rows at \workrowsep 6pt. Measured: the
% guard costs nothing — still 6 pages, and p5 now opens on Part D. Mirrored in the key.
\boxguard[14]
\item Let $g(x) = 2^x - 8$. Give (a) the domain, \ (b) the $y$-intercept, \ (c) the zero, \ (d) the
equation of the horizontal asymptote and the range, \ (e) the end behavior at both ends.
\begin{work}
  \text{(a)} \quad &\text{all real numbers}, \; (-\infty,\infty) \\
  \text{(b)} \quad g(0) &= 2^0 - 8 = 1-8 = -7, \text{ so } (0,-7) \\
  \text{(c)} \quad 2^x - 8 = 0 \;\Rightarrow\; 2^x &= 8 \;\Rightarrow\; x = 3 \\
  \text{(d)} \quad &y = -8; \quad \text{range } (-8,\infty) \\
  \text{(e)} \quad &\text{as } x \to -\infty, \; g \text{ levels off toward } -8;
      \text{ as } x \to \infty, \; g \text{ rises without bound}
\end{work}

\end{enumerate}

% ======================================================================
\parthead{Part D --- Extended Response \quad (2 items, 15 points each --- 30 points)}

\begin{enumerate}[leftmargin=1.8em, itemsep=8pt, topsep=5pt]
\small

\item \textbf{Read one graph completely.} A food truck's daily profit is $P(x)$, in \emph{hundreds
of dollars}, when its signature sandwich is priced at $x$ dollars, for $0 \le x \le 14$.

\begin{center}
\begin{tikzpicture}
\begin{axis}[testaxis, width=11.4cm, height=5.4cm,
    xlabel={$x$ (price, dollars)}, ylabel={$P$ (hundreds of dollars)},
    xmin=0, xmax=14, ymin=-5.4, ymax=7.4,
    xtick={0,2,4,6,8,10,12,14}, ytick={-4,-2,0,2,4,6},
    minor x tick num=1, minor y tick num=1]
  \addplot[royal, very thick] coordinates {(0,-4) (4,0) (7,6) (10,6) (14,-2)};
\end{axis}
\end{tikzpicture}
\end{center}

\begin{enumerate}[label=(\alph*), leftmargin=1.8em, itemsep=4pt, topsep=4pt]
  \item Write the domain and the range in interval notation.
\begin{work}
  \text{domain } [0,14], \qquad \text{range } &[-4,6]
\end{work}

  \item Give both zeros, and say what a zero means for this food truck.
\begin{work}
  x = 4 \text{ and } x &= 13 \\
  \text{these are the two break-even prices} &\text{ --- profit is } \$0 \text{ there}
\end{work}

  \item Give the intervals on which $P$ is increasing, decreasing, and constant.
\begin{work}
  \text{increasing } (0,7), \quad \text{constant } (7,10), \quad \text{decreasing } &(10,14)
\end{work}

  \item Give the absolute maximum as a \textbf{value and a location}. Read the location carefully
        --- it is not a single price.
\begin{work}
  \text{maximum profit } \$600 \; (P=6), \text{ at every price from } &\$7 \text{ to } \$10 \\
  \text{the location is the \emph{interval} } [7,10], \text{ not one point}&
\end{work}

  \item A consultant tells the owner: ``charge as much as you can --- more money per sandwich
        means more profit.'' Use the graph to explain why that is wrong, and name the price you
        would recommend and why.

\writelines{3}
\end{enumerate}

% BOXGUARD: without this, page 6 held nothing but D2(c) and D2(d) — about 1.5in
% of content on an otherwise blank page. That is a boxguard violation, not a page
% count: `make check` was clean at 6/6 throughout, because the key stranded the
% same two parts. A \workrowsep sweep was run first, per Unit 1's rule that on a
% test you reach for \workrowsep before a guard: 5, 4 and 3pt all still gave 6
% pages and only 2pt bought the page back — far too little handwriting room for a
% test whose Part C is ten computations. So the room stays at 6pt and this guard
% sends item 2 over whole instead; p6 now carries the entire item (~5in) rather
% than a stub. Mirrored byte-for-byte in the key. The guard sits in an enumerate,
% not inside a breakable tcolorbox, so \Needspace works normally here.
\boxguard[26]
\item \textbf{Choose a model, use it, and know where it stops.} A fundraiser's post is shared
online. The table gives the number of shares on each of the first four days.

\vspace{3pt}
\begin{center}
\begin{tabular}{@{}l|cccc@{}}
Day $d$ & 0 & 1 & 2 & 3 \\ \hline
Shares $S$ & 8 & 24 & 72 & 216 \\
\end{tabular}
\end{center}

\begin{enumerate}[label=(\alph*), leftmargin=1.8em, itemsep=4pt, topsep=4pt]
  \item Compute the first differences, the second differences, and the ratios. Name the family and
        justify your choice with the row that is constant.
\begin{work}
  \text{first differences: } 16,\; 48,\; 144 &\quad \text{--- not constant} \\
  \text{second differences: } 32,\; 96 &\quad \text{--- not constant} \\
  \text{ratios: } 3,\; 3,\; 3 &\quad \text{--- constant} \;\Rightarrow\; \textbf{exponential}
\end{work}

  \item Write the model $S(d)$, then use it to predict the number of shares on day $5$.
\begin{work}
  S(d) &= 8 \cdot 3^{\,d} \\
  S(5) &= 8 \cdot 3^5 = 8 \cdot 243 = 1944 \text{ shares}
\end{work}

  \item A student says: ``the jumps keep getting bigger --- $16$, then $48$, then $144$ --- so this
        is quadratic.'' Every number the student wrote is correct. Explain what is still wrong with
        the conclusion.

\writelines{3}

  \item This model predicts $S(15) = 8 \cdot 3^{15} = 114{,}791{,}256$ shares --- more people than
        live in most countries. Does that mean the model is \emph{wrong}? Answer in one or two
        sentences using the word \textbf{extrapolation}.

\writelines{2}
\end{enumerate}

\end{enumerate}

\end{document}
